% Transcription — fonds Alexandre Grothendieck, shelfmark 156-6, batch 3 (pages 41-60).
% Established from the facsimile archives/batches/156-6/batch-03.pdf (a copy of
% 156-6.pdf fetched through the project's relay; no cover sheet in the batch,
% so PDF page k = folder page 40+k), per the skill transcribe-grothendieck. The
% folder is ninety-three pages in five batches; this is the third. Batches 1,
% 2, 4 and 5 of this folder and batches 1-5 of the first draft (folder 156-4,
% GF IV) were read for notation when this batch was assembled; the notation of
% batch 2 is kept: 𝔉, 𝓜, ℒ, 𝔓 (with 𝔓_f for closed parts), Φ for a figure as a
% family of parts, compatibility \between, disjointness \parallel, ≪,
% interior refinement \overset{\circ}{\ll}, the underlined triangle
% \trianglelefteq, « omb », « multomb », « Fig », « Ssfig », « Sousat » in roman.
%
% Pass: Opus 5.5 (claude-opus-5-5), 2026-09-26 — first pass, unchecked against the pages by a human.
%
% Pages skipped: none. Pages summarised: none. All twenty leaves are his
% mathematics, in ink, fast drafting, with heavy overwriting on pages 42, 47,
% 49 and 59.
%
% His pagination 41-60 stands at the head of each leaf and coincides with the
% archivists' (noted on page 41, where a struck number precedes the 41). No
% date in his hand on these pages. Numbered runs, recorded where they stand:
% his formula numbers (1.50)-(1.53) (pp. 44-46), (1.54) (cited from p. 48 on
% for the bijection Ssfig(F') → Ssfig(F); the number itself is not written
% beside the formula of p. 48), (1.55)-(1.56) (p. 48), (1.57) (p. 58), (1.58)
% (p. 59); nothing between (1.44), where batch 2 stops, and (1.50) was found
% on these leaves; an unnumbered starred property on p. 42 and a (*) on
% pp. 49, 51 and 60. Axioms: At 1 - At 3 recalled (p. 41), At 4 - At 7
% (p. 42), At 8 (p. 45, « l'axiome des lieux »), At 9 named (p. 59); Sousat 1,
% Sousat 2 (p. 42); At ens 1 (p. 52); Atens*1 - Atens*3 (pp. 54-55) and
% Atens♯1 - Atens♯3 (p. 55, written with a sharp over the asterisk; set
% $^{\sharp}$), « Atens » following batch 4. Internal cross-references,
% as plain page numbers: pages 21, 27 (p. 42), p. 42 (pp. 51, 55), p. 44
% (p. 49), p. 46 (p. 52), p. 47 (p. 50), p. 34 (p. 53), and « cf GF VI
% p. 12 » in the margin of p. 52 (his own chapter number; transcribed as
% written).
%
% What the batch is about. Pages 41-44: sub-ateliers. A bijection between
% « spécieux » closed parts 𝔉' of an atelier 𝔉 and closed parts 𝓜' of 𝓜
% (𝔉' ↦ 𝔉' ∩ 𝓜, 𝓜' ↦ Env(𝓜')); At 1 - At 3 recalled; a sub-atelier is a
% non-empty closed part satisfying Sousat 2 (a condition for F ∨ G to stay in
% 𝔉'), At 4 - At 7 passing trivially; spécieux sub-ateliers; 𝔉 → Fig(𝓜) as
% an example; transitivity (1.50), (1.51); the « trivial » ateliers
% (𝓜, ≤, ≪ = ≤). Pages 45-46: At 8, the axiom of lieux (1.52)
% omb(X)° ≠ ∅; multomb(F) indexed by all of F̃; a corollary (X ≤ Y iff
% omb X ⊂ omb Y iff ...), the bijection (1.53) Ssfig(F) ≅ Ssfig(multomb F);
% « plein » and « ℒ-plein » sub-ateliers. Pages 47-51: quasi-morphisms of
% ateliers (increasing, commuting with bounded Sup and Inf), the condition
% (1.54) that Ssfig(F') → Ssfig(f(F')) be bijective, the map of ordered sets
% (1.55) φ: F̃ → F̃' and (1.56); a lemma on Hom_croiss(I, I') and maps of
% closed parts (p. 49, much overwritten); the trivial ateliers show (1.54)
% is substantial; examples f*: Fig(ℒ') → Fig(ℒ) and 𝔉 → Fig(𝓜) → Fig(ℒ);
% plongements, conditions a'), b'). Pages 52-54: « ensembliste » ateliers
% (At ens 1), the example of topological simplexes (p. 53), and a Scholie:
% a (quasi-)ensemblist atelier is a pair (ℒ, 𝔉), 𝔉 a sub-atelier of Fig(ℒ)
% with {{x}} ∈ 𝔉 (Atens*1), ℒ recovered as the ≪-minimal elements. Pages
% 55-58: Atens*2, *3, the description by (ℒ, 𝓜, ≬_𝓜) with Atens♯1-3,
% Scholie and Corollaire for ateliers of finite type, « finiment spécieux »,
% and the remark (1.57) when X ∩ Y is a common multistrate. Pages 59-60:
% « axiomes de disjonction »: F ∥ G checked on strata, (1.58) (stability of
% ∥ under ≪) to be derived from At 9; an ensemblist counterexample (three
% one-stratum multistrates A, B, A' ⊊ A) shows it does not follow from
% Atens♯3, whence the condition (*) that batch 4 takes up on page 61.
%
% Links with the neighbouring batches. Page 60 ends on the condition (*),
% which page 61 (batch 4) takes up. Batch 4 (p. 72) cites « omb(X)° »
% with the number « (1.48) »: the formula is the one numbered (1.52) here,
% whose label reads fairly clearly « 1.52 ». Batch 5 (p. 91) cites « Sous at
% 1, 2 » on « p. 47 » (second digit doubtful there): Sousat 1 and Sousat 2
% stand on page 42 of this batch (and p. 51 refers to them as « p. 42 »).
%
% Where the second draft restates the first (GF IV, folder 156-4): notes on
% p. 45 (At 8 turns the unproved Proposition of GF IV p. 88 into an axiom),
% p. 54 (lieux as ≪-minimal elements, GF IV p. 87) and p. 59 (disjunction
% axioms, announced at GF IV p. 86).
%
% Reading hazards fixed once here. His looped capital for a figure as a
% family of parts is set \Phi, as in batch 2 (it could be read as a script
% A). « Ssfig » and « Sousfig » both occur and are kept. The overlined ≬ and
% ∥ (p. 60) are negations, set \overline. Slanted left-margin notes on pages
% 41, 42, 46, 47, 49, 50, 51, 52, 57 and 59 were turned and read only in part.
%
% Readings flagged and not settled: the (*) label on p. 42; « sans rien »
% (p. 42); which 𝔉 is primed in the first sentence of p. 44; the top of p. 48
% and its link to (1.54); most of page 49; the opening of page 60
% (« Ça n'a »); « pourtant » (p. 57); « fig. » (p. 56). Variances on the page
% and not repaired: Ssfig(F') written on both sides of the map of p. 48;
% « G ≪ F » for G ∈ 𝔉' in the definition of « plein » (p. 46); F_A ≬ G_A for
% G_B (p. 55); « sous-figure de ℒ » (p. 57).
\documentclass[11pt,a4paper]{article}
\input{../preamble/grothendieck}

\folder{156-6}
\batch{3}
\pages{41}{60}
\dating{1986}
\watermark{Édition de démonstration}
\foldertitle{[Chapitre] VI. Analysis situs (deuxième mouture) : notes manuscrites (18-20/06/1986)}

\begin{document}
\page{41}
\note{pagination de sa main 41 à 60 en tête des feuillets, qui coïncide avec celle des archivistes ; ici, devant le 41, un numéro biffé}
\struck{le sous-atelier $\mathfrak{F}'$ \ill{} partie de $\mathfrak{F}$, comme l'ens. des $\mathfrak{F}$ : $\mathfrak{F}' = \lbrace F \in \mathfrak{F} \mid \widetilde{F} \subset \mathcal{M}' \rbrace \overset{\mathrm{df}}{=} \mathrm{Env}(\mathcal{M}')$. Donc} On trouve ainsi\note{les quatre premières lignes de la page sont barrées de deux longs traits obliques}

\textbf{Proposition} (Ne dépend pas en fait des données $\ll$ sur $\mathfrak{F}$, \add{seulement de $\leq$}) On a deux applications \struck{\ill{}} biunivoques inverses l'une de l'autre
\[
  \begin{array}{ccc}
    & \mathfrak{F}' \longmapsto \mathfrak{F}' \cap \mathcal{M} & \\
    \text{Sousatspéc}(\mathfrak{F}) & \longrightarrow & \mathrm{Ferm}(\mathcal{M}) \\
    & \longleftarrow & \\
    \lbrace F \in \mathfrak{F} \mid \widetilde{F} \subset \mathcal{M}' \rbrace \overset{\mathrm{df}}{=} \mathrm{Env}(\mathcal{M}') & \longleftarrow & \mathcal{M}'
  \end{array}
\]
\note{« Sousatspéc » est écrit sous un « Satsp » biffé ; la flèche du bas, de $\mathcal{M}'$ vers $\mathrm{Env}(\mathcal{M}')$, porte un talon ($\mapsto$)}
entre l'ensemble des \struck{sous-ateliers} \add{parties fermées} \add{spécieuses} $\mathfrak{F}'$ de $\mathfrak{F}$, et l'ens. des parties \emph{fermées} $\mathcal{M}'$ de $\mathcal{M}$.

\marginal{\struck{Dans} \ill{} on identifie \ill{} $\mathfrak{F}$ \ill{} $\mathfrak{P}(\mathcal{M})$, \ill{} \ill{} $\mathrm{Env}(\mathcal{M}') = \lbrace F \in \mathfrak{F} \mid F \subset \mathcal{M}' \rbrace$}\note{note oblique dans la marge gauche, en haut, lue en partie}

Il faut voir, \struck{\ill{}} \add{\ill{} $\mathfrak{F}$ est un atelier, si} $\mathfrak{F}'$ est un sous-atelier de $\mathfrak{F}'$, \struck{donc} pour $\leq$, $\ll$, \struck{c'est un} \add{au sens des} ateliers At 1 -- At 7. Rappelons que ça signifie At 1 -- At 3
\begin{itemize}
  \item[At 1] \struck{\ill{}} et les Sup majorés \add{(pour $\leq$)} existent dans $\mathfrak{F}$ ; et $\mathfrak{F} \neq \emptyset$ (comme on sait, a un plus petit élément $\emptyset_{\mathfrak{F}}$).
  \item[At 2] Toute \struck{fig.} $F \in \mathfrak{F}$ est borne sup. pour $\leq$ des « figures élémentaires » (\ill{} str\supplied{ictement} irr.) qu'elle majore (pour $\leq$)
  \item[At 3] Si $F, G \in \mathfrak{F}$, pour que $F \vee G$ existe dans $\mathfrak{F}$ \add{(i.e. $\lbrace F, G \rbrace$ majoré)}, il suffit que \struck{\ill{}} pour $X \trianglelefteq F$, $Y \trianglelefteq G$, $X \vee Y$ existe.
\end{itemize}

\page{42}
Pour At 1, At 2 c'est clair, pour At 3 il faut \struck{de plus} imposer, en plus de

\begin{itemize}
  \item[Sousat 1] $\mathfrak{F}'$ est une partie fermée \add{non vide} de $\mathfrak{F}$,
\end{itemize}
la condition
\begin{itemize}
  \item[Sousat 2] Si $F, G \in \mathfrak{F}'$ et si pour \struck{tout} $X \in \widetilde{F}$, $Y \in \widetilde{G}$, $X \vee Y$ existe dans $\mathfrak{F}'$ (i.e. \struck{$X \between Y$ dans $\mathfrak{F}$, dans $X \vee Y$ existe dans, $X \vee Y \in \mathfrak{F}'$ ; $X \between Y \notin \mathfrak{F}'$, alors} $X \between Y$ dans $\mathfrak{F}$ et \struck{$X \vee Y \in \mathfrak{F}'$} $X$, $Y$ majorés dans $\mathfrak{F}'$) alors $F \vee G \in \mathfrak{F}'$ (NB on sait que $F \vee G$ existe dans $\mathfrak{F}$) (cf. pages 21, 27).
\end{itemize}
\note{au milieu de Sousat 2, deux lignes et demie sont biffées de traits horizontaux, et un petit cadre tracé à gauche est barré de traits obliques ; ce qui se lit est donné comme biffé}

\marginal{NB (1-2) une condition \ill{} ainsi \ill{} que $(\mathfrak{F}', \leq)$ \ill{} \ill{} condition \ill{} At 3 pour $\mathfrak{F}'$}\note{note oblique dans la marge gauche, en face de Sousat 2, lue en partie}

Les conditions At 4, At 5, At 6, At 7 passent \add{\uncertain{sans rien}} trivialement \struck{de} $\mathfrak{F}$ à $\mathfrak{F}'$. Pour At 7, je note seulement une propriété
\begin{itemize}
  \item[\uncertain{(*)}] si $X, Y \in \mathcal{M}'$, alors $X \overset{\circ}{\ll} Y$ dans $\mathfrak{F}$ équivaut à $X \overset{\circ}{\ll} Y$ dans $\mathfrak{F}'$.
\end{itemize}

\struck{Supposons maintenant $\mathfrak{F}$ interprété comme un ensemble de figures (en fait, un sous-atelier) de $\mathrm{Fig}(\mathcal{M})$.}\note{ces trois lignes sont barrées de traits obliques}

\textbf{Définition} Soit $\mathfrak{F}$ un atelier au sens At 1 -- At 7. On \struck{appelle} \emph{sous-atelier} de $\mathfrak{F}$, une partie $\mathfrak{F}'$ de $\mathfrak{F}$ \struck{\ill{}}, satisfaisant Sousat 1 et Sousat 2 ci-dessus, i.e. telle que 1°) $\mathfrak{F}'$ soit une partie fermée non vide de $(\mathfrak{F}, \leq)$ \struck{et} (ce qui s'exprime par Sousat 1) et 2°) $\mathfrak{F}'$, muni de $\leq$, $\ll$, est un atelier (ce qui équivaut : Sousat 2). On dit que

\page{43}
$\mathfrak{F}'$ est un \emph{sous-atelier} \emph{spécieux}, si c'est une partie fermée \emph{spécieuse} de $\mathfrak{F}$ \struck{Conséquence la condition Sousat 2 est} Ce qui suffit à entraîner que c'est un sous-atelier.

\textbf{Exemple} Soit $\mathfrak{F}$ un atelier, et considérons l'application canonique
\[
  \mathrm{Multomb} : \mathfrak{F} \longrightarrow \mathrm{Fig}(\mathcal{M}),
\]
on a vu que l'image est une partie fermée pour $\leq$, et que les relations $\leq$ et $\ll$ sur $\mathfrak{F}$ sont induites par celles de $\mathrm{Fig}(\mathcal{M})$. Ainsi $\mathfrak{F}$ est un sous-atelier de l'atelier des figures \struck{dans} \add{dans} $\mathcal{M}$. \add{Mais ce n'est pas en général} \struck{En fait, on a vu aussi que c'est un} sous-atelier spécieux (mais ça l'est peut-être dans le cas ensembliste, à voir).

\struck{Je reviens à la \ill{} aux « petits » \ill{}}

\textbf{Remarque} Soit $\mathfrak{F}$ un ordonné par $\leq$, $\mathfrak{F}'$ une partie fermée \add{non vide} de $\mathfrak{F}$. Alors les parties fermées non vides de $\mathfrak{F}'$ sont exactement les parties fermées non vides de $\mathfrak{F}$ qui sont contenues dans $\mathfrak{F}'$. On a le même énoncé pour les parties fermées spécieuses, pourvu que l'on suppose $\mathfrak{F}'$ spécieux. De façon précise, si \ill{} on a des

\page{44}
inclusions,
\[
  \mathfrak{F}'' \subset \mathfrak{F}' \subset \mathfrak{F}
\]
de parties fermées non vides, si $\mathfrak{F}''$ est spécieux dans \uncertain{$\mathfrak{F}$}, il l'est dans $\mathfrak{F}'$,\note{les deux $\mathfrak{F}$ de cette proposition sont mal distingués de $\mathfrak{F}'$ sur la page ; la lecture suit le sens (la réciproque demande $\mathfrak{F}'$ spécieux)} et la réciproque est vraie si $\mathfrak{F}'$ est spécieux dans $\mathfrak{F}$. De même, si $\mathfrak{F}$ est un atelier, $\mathfrak{F}'$ un sous-atelier, les sous-ateliers de $\mathfrak{F}'$ sont les sous-ateliers de $\mathfrak{F}$ contenus dans $\mathfrak{F}'$ --- en d'autres termes, la condition supplémentaire Sousat 2 pour $\mathfrak{F}''$ est la même, qu'on la prenne par rapport à $\mathfrak{F}'' \subset \mathfrak{F}'$ ou par $\mathfrak{F}'' \subset \mathfrak{F}$ ; en effet, c'est une condition intrinsèque : $\mathfrak{F}''$ \ill{} (\uncertain{savoir}, At 3 pour $\mathfrak{F}''$). Donc, \struck{\ill{}}
\begin{align*}
  &(1.50) & \mathrm{Sousat}(\mathfrak{F}') &= \lbrace \mathfrak{F}'' \in \mathrm{Sousat}(\mathfrak{F}) \mid \mathfrak{F}'' \subset \mathfrak{F}' \rbrace
\end{align*}
et si $\mathfrak{F}'$ spécieux dans $\mathfrak{F}$
\begin{align*}
  &(1.51) & \text{Sousatspéc}(\mathfrak{F}') &= \lbrace \mathfrak{F}'' \in \text{Sousatspéc}(\mathfrak{F}) \mid \mathfrak{F}'' \subset \mathfrak{F}' \rbrace .
\end{align*}
\note{numéros de formules de sa main, écrits « (1,50) » et « (1;51) »}

\textbf{Remarque} Soit $(\mathfrak{F}, \leq)$ un ordonné satisfaisant At 1 -- At 3, \struck{donc} qui se décrit donc par $(\mathcal{M}, \leq, \mathfrak{F})$, \struck{$\mathfrak{F}$ \ill{}} $(\mathcal{M}, \leq)$ un ensemble ordonné ; $\mathfrak{F} \subset \mathfrak{P}_{\mathrm{f}}(\mathcal{M})$ \add{ensemble des parties fermées de $\mathcal{M}$} telle que 1°) $|\mathfrak{F}| = \mathcal{M}$ i.e. $\forall x \in \mathcal{M}$, $\exists\, F \in \mathfrak{F}$ avec $x \in F$, et 2°) $\mathfrak{F}$ fermé dans $\mathfrak{P}_{\mathrm{f}}(\mathcal{M})$, i.e. si $F \in \mathfrak{F}$ et $G$ partie fermée de $F$, alors $G \in \mathfrak{F}$. Prenons dans $\mathcal{M}$ la relation $\ll \overset{\mathrm{df}}{=} \leq$. Alors, si $(\mathcal{M}, \leq)$ satisfait At 3, alors $(\mathcal{M}, \leq, \ll)$ satisfait At 1 à At 7. Pour moi

\page{45}
axiomes qui vont suivre, il y en a qui \uncertain{les} éliminent, je pense, cet exemple trivial d'atelier.

\note{un court trait horizontal sépare ce qui précède de ce qui suit}

Revenons aux \emph{petites} ombres et multiombres. C'est le moment d'introduire\note{dans la première mouture (GF IV, dossier 156-4, p. 88), la condition $A^{\circ} = A \smallsetminus \partial A \neq \emptyset$ pour les membres de la multiombre figure dans une proposition énoncée sans démonstration ; elle est posée ici en axiome, At 8, sous la forme (1.52)}
\begin{itemize}
  \item[\emph{At 8}] Soit $X \in \mathcal{M}$, et soit $\partial X = \operatorname*{Sup}^{(\leq)}_{\substack{Y \trianglelefteq X \\ Y \neq X}} Y$. Alors $\exists\, x \in \mathcal{L}$ tel que $x \ll X$, \struck{$x \notin$} $x \not\ll \partial X$ (i.e. $x \overset{\circ}{\ll} X$). \struck{et En}
\end{itemize}
En d'autres termes, posant \struck{\ill{}}
\[
  \begin{aligned}
    \mathrm{omb}(X)^{\circ} &= \mathrm{omb}(X) \smallsetminus \mathrm{omb}(\partial X) \\
    &= \mathrm{omb}(X) \smallsetminus \bigcup_{\substack{Y \trianglelefteq X \\ Y \neq X}} \mathrm{omb}(Y),
  \end{aligned}
\]
on a :
\[
  (1.52) \qquad \boxed{\ \mathrm{omb}(X)^{\circ} \neq \emptyset\ }
\]

Considérons alors, pour $F \in \mathfrak{F}$, \add{la} multiombre, qui est par définition \add{la figure ensembliste} \struck{partie \ill{}}
\[
  \mathrm{multomb}(F) = \lbrace \mathrm{omb}(X) \mid X \trianglelefteq F,\ \mathrm{omb}(X)^{\circ} \neq \emptyset \rbrace .
\]
\note{sous « $X \trianglelefteq F$ » : « i.e. $X \in \widetilde{F}$ »}
At 8 nous dit que l'ensemble d'indices ici est $\widetilde{F}$ tout entier. \add{(cf. note ensembliste)}\note{addition entre les lignes, lue en partie} Ainsi on trouve

\textbf{Corollaire} \struck{Soient} $F$ une figure, $X, Y \in \widetilde{F}$. Alors conditions suivantes équivalentes
\begin{itemize}
  \item[a)] $X \leq Y$ \quad ($\Longleftrightarrow X \ll Y$)
\end{itemize}

\page{46}
\begin{itemize}
  \item[b)] $\mathrm{omb}(X) \subset \mathrm{omb}(Y)$
  \item[c)] $\mathrm{omb}(X)^{\circ} \subset \mathrm{omb}(Y)$
  \item[d)] $\mathrm{omb}(X)^{\circ} \cap \mathrm{omb}(Y) \neq \emptyset$
\end{itemize}
Il y a une correspondance biunivoque
\[
  (1.53) \quad \left\lbrace
  \begin{array}{c}
    F' \longmapsto \mathrm{multomb}(F') \\
    \mathrm{Ssfig}(F) \xrightarrow{\ \sim\ } \mathrm{Ssfig}(\mathrm{multomb}(F))
  \end{array}
  \right.
\]
\note{sous « $\mathrm{Ssfig}(F)$ », un signe $=$ vertical et « $\mathfrak{F}_{\leq F}$ »}

\textbf{Corollaire} Pour que $F$ soit élémentaire, i.e. une multistrate, il faut et il suffit que $\mathrm{multomb}(F)$ le soit.

\textbf{Remarque} Je n'ai aucune raison que At 8 soit stable par passage à un sous-atelier $\mathfrak{F}'$, \struck{\ill{}} pourvu que celui-ci soit spécieux. Mais \struck{considérons} supposons que \struck{$G$} $F \in \mathfrak{F}'$, $G \in \mathfrak{F}$, $G \ll F$ implique $G \ll F$\note{sic ; on attend « implique $G \in \mathfrak{F}'$ »} i.e. que $\mathfrak{F}'$ soit fermé pour $\ll$ également --- on dit alors que le sous-atelier est \emph{plein}. En général, on aura
\[
  \mathcal{L} \cap \mathfrak{F}' \subset \mathcal{L}' \quad (\overset{\mathrm{df}}{=} \text{ens. des lieux de } \mathfrak{F}')
\]
mais dans le cas d'un atelier plein,$^{*}$ on aura
\[
  \mathcal{L} \cap \mathfrak{F}' = \mathcal{L}' .
\]
Dans ce cas,$^{*}$ si $\mathfrak{F}$ satisfait At 8, $\mathfrak{F}'$ aussi. De plus, \struck{---} pour $F \in \mathfrak{F}'$, $\mathrm{omb}(F)$ et $\mathrm{multomb}(F)$ sont les mêmes, au sens de $\mathfrak{F}$ ou au sens de $\mathfrak{F}'$.

\marginal{$^{*}$ il suffit qu'il soit $\mathcal{L}$-plein, i.e. $F \in \mathfrak{F}'$, $x \in \mathcal{L}$, $x \ll F \Rightarrow x \in \mathfrak{F}'$}\note{note oblique dans la marge gauche ; les deux astérisques du texte y renvoient}

\page{47}
\struck{\textbf{Corollaire 2} Soient $F, G \in \mathfrak{F}$, \ill{} $G < F$. Alors $\exists\, x \in \mathrm{omb}(F)$, \ill{}}

\struck{Je reviens à la notion d'atelier \emph{ensembliste}}\note{ces lignes, en haut de la page, sont biffées horizontalement et barrées de traits obliques, ainsi qu'un « Déf » commencé en marge}

Je reprends la notion des \struck{homomorphismes} morphismes d'ateliers.

\textbf{Définition} Soient $\mathfrak{F}'$, $\mathfrak{F}$ deux ateliers. On appelle \struck{\ill{}} \emph{quasi-morphisme} \struck{de} (car. \ill{} \ill{} un morphisme !) de $\mathfrak{F}'$ dans $\mathfrak{F}$, une application
\[
  f : \mathfrak{F}' \longrightarrow \mathfrak{F}
\]
qui est
\begin{itemize}
  \item[a)] croissante pour $\leq$ et $\ll$
  \item[b)] commute aux Sup \struck{majorés} (majorés) (et \add{donc} transforme $\emptyset_{\mathfrak{F}'}$ en $\emptyset_{\mathfrak{F}}$) et aux Inf pour $\leq$.
\end{itemize}

\marginal{il vaudrait mieux de parler de « quasi-morphismes » (en réservant le nom de morphismes \ill{} plus bas)}\note{note oblique dans la marge gauche, en haut}

On n'exige pas qu'elle transforme fig. élémentaires en fig. élémentaires. Mais quand il en est ainsi, cette application est connue quand on connaît sa restriction \struck{restriction}
\[
  f_{\mathrm{str}} : \mathcal{M}' \longrightarrow \mathcal{M}
\]
aux ens. de multistrates \add{(= fig. élém.)}. Dans ce cas, pour que $f$ soit injective, il faut et il suffit que $f_{\mathrm{str}}$ le soit (et en termes de plongement \ill{}) \ill{} \ill{} : $\mathfrak{F} \hookrightarrow \mathfrak{P}(\mathcal{M})$, $\mathfrak{F}' \hookrightarrow \mathfrak{P}(\mathcal{M}')$, $f$ s'identifie alors à la restriction \struck{\ill{}} de $f : A \mapsto f_{\mathrm{str}}(A) : \mathfrak{P}(\mathcal{M}') \to \mathfrak{P}(\mathcal{M})$, \add{du moins quand $f(\mathcal{M}')$ \ill{} une partie} \add{\ill{} fermée de $\mathcal{M}$, cf plus bas}. \ill{} \ill{} \ill{} donne ensuite \ill{} \ill{} \ill{} quasi-morphl. des conditions simples, qui indiquent que lorsque $f$ est injectif, alors l'application induite :

\marginal{il pas clair \ill{} faudrait $F' \in \mathfrak{F}'$, $f(F') = \ldots$ \ill{} voir plus bas}\note{note oblique dans la marge gauche, en face des dernières lignes, lue en partie ; un trait vertical la sépare du texte}

\page{48}
\[
  \begin{array}{ccc}
    \mathrm{Ssfig}(F') & \longrightarrow & \mathrm{Ssfig}(F') \\
    G' & \longmapsto & f(G')
  \end{array}
  \qquad \text{où } F = f(F')
\]
\note{il écrit $\mathrm{Ssfig}(F')$ des deux côtés ; le second est sans doute $\mathrm{Ssfig}(F)$, comme l'indique « où $F = f(F')$ ». La marge renvoie à cette application sous le numéro (1.54), qui n'est pas écrit ici}
\emph{bijective}. Sans supposer $f$ injective, dans les cas à ma connaissance, on aura une application canonique (croissante)

\begin{tikzcd}
  \widetilde{F} \arrow[d, "{\varphi_{f,F'} = \varphi}"] \\
  \widetilde{F}'
\end{tikzcd}
\note{numéro (1.55) en marge du diagramme ; devant $\widetilde{F}$, un signe souligné, peut-être un chiffre, non lu}

\ill{} \ill{} (1.54) un \ill{} d'un ordonnés, de telle façon, que (1.54) se déduit par l'application en sens inverse
\[
  (1.56) \quad \left\lbrace
  \begin{array}{c}
    \mathfrak{P}_{\mathrm{f}}(\widetilde{F}') \longrightarrow \mathfrak{P}_{\mathrm{f}}(\widetilde{F}) \\
    A \longmapsto \varphi^{-1}(A)
  \end{array}
  \right. ,
\]
et de plus, la \struck{topologie de} relation d'ordre sur $\widetilde{F}$ est image inverse de celle sur $\widetilde{F}'$. Dans le cas où $f$ est injective, et où la (1.54) est satisfaite, (1.55) sera un iso. d'un ordonnés. Dans le cas général, on a vu que l'application (1.54) est \struck{\ill{}} au moins \emph{surjective}. \struck{\&} Mais dans \struck{tout} un \ill{} (comme $\mathfrak{F} \to \mathrm{Fig}(\mathcal{L})$), des cas \uncertain{importants} même si $f$ n'est pas injective, on aura (1.54) bijective pour tout $F$.

On notera que dans les cas où (1.54) peut se déduire d'un \struck{\ill{}} morphisme d'ordonnés (1.55), celui-ci est unique. En d'autres termes

\marginal{\struck{\ill{}} \ill{} partie (\ill{}) \ill{} non \ill{}, \ill{} que (1.54) \ill{} pourtant}\note{note oblique dans la marge gauche, en haut, lue par fragments}

\page{49}
\textbf{Lemme} Soient $I$, $I'$ deux ens. \struck{pré}ordonnés. Alors l'application
\[
  \mathrm{Hom}_{\mathrm{croiss}}(I, I') \longrightarrow \mathrm{Hom}_{\mathrm{cr}}(\mathfrak{P}_{\mathrm{f}}(I'), \mathfrak{P}_{\mathrm{f}}(I))
\]
entre applications croissantes de $I$ dans $I'$, et applic. cr. de $\mathfrak{P}_{\mathrm{f}}(I)$ en $\mathfrak{P}_{\mathrm{f}}(I')$ (parties fermées de $I$, $I'$) est \emph{injective}. Son image \struck{\ill{}} est formée des \add{contenue dans les}
\[
  (*) \quad \mathfrak{P}_{\mathrm{f}}(I') \longrightarrow \mathfrak{P}_{\mathrm{f}}(I)
\]
\add{et réunions} commutant aux intersections \add{(quelconques)} \struck{(\ill{}, \ill{} kif-kif)}, \struck{et réunions \ill{}} \add{\ill{}} si $I$ est sobre (p.ex. $I$ fini) on voit par \ill{} \ill{}.
\note{la page est très retravaillée : lignes biffées horizontalement, additions interlinéaires et notes marginales ; les indications « de $\mathfrak{P}_{\mathrm{f}}(I)$ en $\mathfrak{P}_{\mathrm{f}}(I')$ » et la flèche (*) vont en sens contraires, comme sur la page}

\struck{\ill{}} Preuve \ill{} \ill{} complémentaire \ill{}, \add{\ill{} \ill{} des} \add{\ill{}} applications (*) correspondant aux \struck{quasi-compatibilités \ill{} \ill{}} \add{\ill{} \ill{}} applications
\[
  \mathrm{Top}(I) \longrightarrow \mathrm{Top}(I')
\]
de Topos \struck{\ill{} (compatibilité aux \ill{} \ill{} i.e. sobre)}, \add{intersections quelconques} \struck{$I$ est ordonné, i.e. sobre} applications continues \add{i.e. croissantes} \ill{} \ill{} $I \to \widehat{I}'$ ($\widehat{I}'$ : $I'$ sobrifié) \ill{} supposons $I$ sobre) \struck{$\widetilde{I}$}.

\struck{NB on n'avait pas \ill{} $\to \widetilde{I}'$ ($I$ sobre) \ill{} \ill{} \ill{} \ill{} i.e. \ill{} ; en effet pour ce \uncertain{bijectivité} \ill{}. Donc cela nous fournit bien (1.55) p.ex. $I'$ fini, au moins si $I'$ sobre}\note{ce paragraphe est biffé de plusieurs traits horizontaux et lu par fragments}

Mais aura-t-on nécessairement que la top de $\widetilde{F}$ est image inv. de celle de $\widetilde{F}'$ ? La surjectivité \ill{}. En fait, c'est \emph{équivalent} à la surjectivité de (1.54).

\marginal{Elles commutent aux réunions quelconques, et \ill{} \ill{} \ill{} ; \ill{} \ill{} \ill{} des parties \ill{} \ill{} forme $\widehat{I}'$ \ill{} $= I'_{\mathrm{sobr}}$ \ill{} \ill{}}\note{deux notes obliques dans la marge gauche, l'une en haut, l'autre en face du NB biffé, cette dernière reliée au texte par une accolade ; lues par fragments}

\textbf{Ex} Prenons des ateliers « triviaux » (p. 44)
\[
  (\mathcal{M}, \leq, \ll \overset{\mathrm{df}}{=} \leq, \mathfrak{F}).
\]
\struck{\ill{} \ill{}} définis par \struck{\ill{}} $\mathfrak{F}$ : l'axiome At 8, il signifie que pour \struck{tel} $\forall\, X \in \mathcal{M}$, $\exists\, x \in \mathcal{M}$, $x$ \emph{minimal}, tel que $x \leq X$, mais et que \struck{$X$} soit « successeur » de \add{(dans le cas $X \neq x$ i.e. $X$ non minimal)}

\page{50}
de $X$. i.e. $\nexists\, Y \in \mathcal{M}$ avec $x \leq Y < X$. C'est évidemment une condition des plus restrictives sur $\mathcal{M}$ ! Par exemple \ill{} \ill{} \ill{} \ill{} de At 8. Considérons deux \struck{tels} ateliers
\[
  (\mathcal{M}', \leq, \ll \overset{\mathrm{df}}{=} \leq, \mathfrak{F}'),
\]
prenons \struck{tous} p.ex. pour $\mathfrak{F} = \mathfrak{P}_{\mathrm{f}}(\mathcal{M})$, $\mathfrak{F}' = \mathfrak{P}_{\mathrm{f}}(\mathcal{M}')$, et intéressons-nous aux
\[
  f : \mathfrak{F}' \longrightarrow \mathfrak{F}
\]
\struck{tels} qui soient des morphismes d'ateliers, au sens de la définition p. 47. Par le lemme, ils correspondent aux \struck{\ill{}} applications croissantes en sens inverse

\begin{tikzcd}
  \mathcal{M} \arrow[d, "\varphi"] \\
  \mathcal{M}'
\end{tikzcd}

sans aucune condition supplémentaire. Cela montre que la condition que \ill{} \ill{} la top. de $\varphi$ dans (1.55) soit telle que la top. de $\widetilde{F}$ soit l'image inverse pour $\varphi$, i.e. que (1.54) soit toujours surjective, est bien substantielle.

\marginal{\ill{} \ill{} de $\mathcal{M}$ \ill{} \ill{} \ill{}}\note{courte note oblique dans la marge gauche, en haut, presque illisible}

\textbf{Exemples} 1) Soit $f : \mathcal{L} \to \mathcal{L}'$ application ensembliste, elle induit \struck{Fig}
\[
  \begin{array}{rcl}
    f^{*} : \mathrm{Fig}(\mathcal{L}') & \longrightarrow & \mathrm{Fig}(\mathcal{L}) \\
    \Phi' & \longmapsto & \lbrace f^{-1}(A') \mid A' \in \Phi',\ f^{-1}(A') \neq \emptyset \rbrace
  \end{array}
\]
\note{sous le $\Phi'$ de la seconde ligne, surchargé et souligné, un signe à peine lisible}

\page{51}
qui est un quasi-morphisme d'ateliers

2) $\boxed{\mathfrak{F} \to \mathrm{Fig}(\mathcal{M})}$ est un quasi-morphisme d'ateliers injectif, satisfaisant la condition \add{de bijectivité} (1.54). \add{\ill{}} En ce composant avec le quasi-morphisme Fig
\[
  \mathrm{Fig}(\mathcal{M}) \longrightarrow \mathrm{Fig}(\mathcal{L})
\]
induit comme dans 1° de $\mathcal{L} \hookrightarrow \mathcal{M}$, on trouve le quasi-morphisme
\[
  \mathfrak{F} \longrightarrow \mathrm{Fig}(\mathcal{L})
\]
\struck{l'application} \add{il} satisfait la condition \add{de bijectivité} (1.54), si cette fois-ci on suppose \emph{At 8}, et \ill{} \ill{} que $f$ est aussi injectif.

\textbf{Définition} On dit qu'un quasi-morphisme $f : \mathfrak{F}' \to \mathfrak{F}$ d'ateliers est un \emph{plongement}, si
\begin{itemize}
  \item[] \struck{a) $f$ est injectif}
  \item[a)] l'image $f(\mathfrak{F}')$ est un sous-atelier de $\mathfrak{F}$ (p. 42)
  \item[b)] $f$ induit un iso. $\mathfrak{F}' \xrightarrow{\ \sim\ } f(\mathfrak{F}')$.
\end{itemize}
\struck{Les} Les conditions \struck{a), b)} équivalent \add{\ill{}} aux \ill{} \ill{} : \add{posant $F = f(F')$}
\begin{itemize}
  \item[a')] $\forall F' \in \mathfrak{F}'$, l'application
  \[
    (*) \quad \mathrm{Sousfig}(F') \longrightarrow \mathrm{Sousfig}(F), \qquad G' \longmapsto f(G')
  \]
  est bijective \add{(en fait, vu b') la surjectivité suffit)} \add{\ill{} \ill{} a' \ill{} \ill{} par l'image}, i.e. $\forall G \leq F$, $\exists\, G' \leq F'$ telle que $G = f(G')$.
  \item[b')] $\forall F', G' \in \mathfrak{F}'$, on a
  \[
    f(G') \ll f(F') \Longrightarrow G' \ll F'
  \]
\end{itemize}

\marginal{\uncertain{Convient-il} \ill{} \ill{} p. 47 ?}\marginal{NB Si $f$ injectif, a) \ill{} a') $\Rightarrow$ \ill{} plus bas.}\note{deux notes obliques dans la marge gauche, lues en partie}

\page{52}
C'est évidemment nécessaire. Pour ce suffisant, \struck{\ill{}} il suffit de voir que \add{$f$ injectif (de façon visible) de b') et que} a), b') impliquent, pour $F', G' \in \mathfrak{F}'$
\[
  f(G') \leq f(F') \Longrightarrow G' \leq F'
\]
\note{« Il suffit de voir que » est encadré ; l'addition interlinéaire est elle-même encadrée en partie}

\emph{Dém} \struck{Par} Par b'), \struck{\ill{}} si $f(G') \leq f(F')$, $\exists\, \overline{G'} \leq F'$ tel que $f(\overline{G'}) = f(G')$. Mais alors par b) on a $G' = \overline{G'}$, d'où $G' \leq F'$, cqfd.

\textbf{Définition} On dit que l'atelier $\mathfrak{F}$ \add{(satisfaisant At 1 -- At 8)} est \struck{\ill{} \ill{}} \emph{ensembliste} (ou mieux : \ill{} « ensembliste » plus \ill{}) si le quasi-morphisme canonique
\[
  \mathfrak{F} \longrightarrow \mathrm{Fig}(\mathcal{L})
\]
est un \emph{plongement}, ou, ce qui revient au même (compte tenu de (1.53) p. 46) si on a
\begin{itemize}
  \item[\emph{At ens 1}] $\forall F, G \in \mathfrak{F}$,
  \[
    \mathrm{multomb}(F) \ll \mathrm{multomb}(G) \Longrightarrow F \ll G .
  \]
\end{itemize}

\marginal{Pour \ill{} \ill{} cf GF VI p. 12}\note{note oblique dans la marge gauche, « VI » encadré ; renvoi à sa propre pagination du chapitre VI (ce dossier), transcrit tel quel}

\textbf{Remarques} 1) On voit qu'il suffit de poser At ens 1 pour $F$, $G$ des multistrates \struck{pour} $X, Y \in \mathcal{M}$. Ce serait plus joli de poser At ens \ill{} \ill{} en termes.

2) On n'exige pas que le sous-atelier $\mathfrak{F} \hookrightarrow \mathrm{Fig}(\mathcal{L})$ soit spécieux. \emph{Par exemples} \struck{ex \ill{} \ill{} B} \ill{}\note{la disposition de la remarque 2) est incertaine : « que le sous-atelier » termine la ligne qui précède « 2) On n'exige pas que \ldots\ spécieux » ; devant « Soit $E$ », un signe cerclé non lu} Soit $E$ un espace topologique \ill{} modéré, et $\mathcal{M}$ l'ens. de ses multistrates \emph{simplexes} (topologiques, ou modérés), mais

\page{53}
(contrairement à l'exemple B) p. 34) sans structure affine. On définit la compatibilité des simplexes comme loc. cit. \struck{Mais on n'a pas \ill{} c'est} D'où un ens. de multistrates, i.e. \add{prendre pour} $\mathfrak{F}$ ensemble des \struck{multistrates} \add{simplexes} fermés pour $\leq$, et fermés de simplexes deux à deux compatibles. L'axiome At 8 trivialement \ill{} (NB $E$ s'identifie à $\mathcal{L}$). Mais notons \ill{} la notion de compatibilité des simplexes est plus fine que la compatibilité des figures ensemblistes associées, puisque ici on exige \add{de plus} que la figure intersection soit élémentaire, i.e. soit une multistrate commune. On y reviendra par la suite\ldots

Ainsi, dans le cas d'un atelier ensembliste, $\mathfrak{F}$, $\mathfrak{F}$ se reconstitue à partir de l'ens. des figures dans $\mathcal{L}$ qu'il décrit, comme un sous-atelier de $\mathrm{Fig}(\mathcal{L})$. La question est donc \emph{quels} sous-ateliers d'un $\mathrm{Fig}(\mathcal{L})$ on obtient ainsi.

\page{54}
\textbf{Scholie} Se donner un atelier (quasi-)ensembliste, « revient au même » que de se donner un couple $(\mathcal{L}, \mathfrak{F})$, où $\mathcal{L}$ est un ens., $\mathfrak{F}$ un sous-atelier de $\mathrm{Fig}(\mathcal{L})$, satisfaisant la condition suivante :
\begin{itemize}
  \item[] \struck{Atens*1 \ill{} Les strates des $\Phi \in \mathfrak{F}$ séparent les points de $\mathcal{L}$, i.e. si $x, y \in \mathcal{L}$, $x \neq y$, alors $\exists\, \Phi \in \mathfrak{F}$, $A \in \Phi$, tel que \ill{} $x \in A$, $y \notin A$ ou $y \in A$, $x \notin A$.}
  \item[\emph{Atens*1}] \struck{\ill{}} $\forall x \in \mathcal{L}$, $\lbrace \lbrace x \rbrace \rbrace \in \mathfrak{F}$
\end{itemize}
\note{la première forme de Atens*1 est barrée de plusieurs traits obliques ; la seconde la remplace}

\struck{Il} Il \struck{\ill{}} est à vérifier que dans ce cas, $\mathcal{L}$ se reconstitue à partir de $(\mathfrak{F}, \leq, \ll)$ comme l'ens. des lieux de $\mathfrak{F}$, donc que
\[
  x \longmapsto \lbrace \lbrace x \rbrace \rbrace = \Phi_x
\]
est une bijection de $\mathcal{L}$ \struck{avec l'ens de $\mathcal{M}$} \add{avec} l'ens. des éléments de $\mathfrak{F} \smallsetminus \lbrace \emptyset \rbrace$ minimaux pour $\ll$\note{c'est la définition des lieux de la première mouture (GF IV, dossier 156-4, p. 87) : les éléments de $\mathfrak{F} \smallsetminus \lbrace \emptyset_{\mathfrak{F}} \rbrace$ minimaux pour $\ll$} \add{\ill{} minimaux dans $\mathrm{Fig}(\mathcal{L}) \smallsetminus \lbrace \emptyset \rbrace$}. Il est clair que $\Phi_x$ est minimal. \struck{Or} Soit $F \in \mathfrak{F} \smallsetminus \lbrace \emptyset \rbrace$, \add{\uncertain{minimal} dans $\mathfrak{F} \smallsetminus \lbrace \emptyset \rbrace$}. \struck{Soit $A \in F$, on a $A^{\circ} \neq \emptyset$, soit $x \in A^{\circ}$ alors} On a $|F| \neq \emptyset$ (puisque $F \neq \emptyset$) soit $x \in |F|$, alors $\Phi_x \ll F$, donc $\Phi_x = F$, cqfd.

\page{55}
Je rappelle pour mémoire les conditions \add{Sousat 1, 2} (de la p. 42, pour qu'une partie $\mathfrak{F}$ de $\mathrm{Fig}(\mathcal{L})$ \ill{} l'\ill{}) :
\begin{itemize}
  \item[\emph{Atens*2}] $\mathfrak{F}$ est fermé dans $\mathrm{Fig}(\mathcal{L})$ pour $\leq$, i.e. si $F \in \mathfrak{F}$, alors toute sous-figure de $F$ est dans $\mathfrak{F}$.
  \item[\emph{Atens*3}] Si $F, G \in \mathfrak{F}$, et si pour \struck{tout $F$ et $G$ sont compat} toutes multistrates $F_A$ de $F$, $G_B$ de $G$, $\lbrace F_A, G_B \rbrace$ est majoré dans $\mathfrak{F}$ (i.e. $F_A \between G_A$ et $F_A \cup G_B \in \mathfrak{F}$\note{il écrit $G_A$ puis $G_B$}, \struck{alors} \add{i.e. \ill{}} \ill{}) \add{en figures} $F$, $G$ sont \add{dans} \ill{} ens. multistrates compatibles), \struck{$F$, $G$} $F \vee G \in \mathfrak{F}$.
\end{itemize}

Ceci peut s'expliciter \add{(partiellement)} aussi en termes d'un ensemble $\mathcal{M}$ de figures \add{(ensemblistes)} élémentaires dans $\mathcal{L}$, \struck{\ill{} \ill{} par $\mathfrak{F}$} \add{relation refl. et sym.} (« compatibilité géométrique ») $\underset{\mathcal{M}}{\between}$ dans $\mathcal{M}$, satisfaisant les conditions
\begin{itemize}
  \item[\emph{Atens$^{\sharp}$1}] $\forall x \in \mathcal{L}$, $\lbrace \lbrace x \rbrace \rbrace \in \mathcal{M}$
  \item[\emph{Atens$^{\sharp}$2}] \struck{\ill{}} $\forall X \in \mathcal{M}$, toute sous-multistrate est dans $\mathcal{M}$
  \item[\emph{Atens$^{\sharp}$3}] Si $X, Y \in \mathcal{M}$ et $X \underset{\mathcal{M}}{\between} Y$, et $X'$, \struck{$Y'$} $\leq X$, $Y' \leq Y$, alors $X' \underset{\mathcal{M}}{\between} Y'$ ; de plus $X \underset{\mathcal{M}}{\between} Y \Longrightarrow X \between Y$.
\end{itemize}
\note{il écrit « At ens » suivi d'un dièse, tracé sur l'astérisque des conditions précédentes ; un trait vertical en marge réunit Atens$^{\sharp}$2 et Atens$^{\sharp}$3}

\page{56}
On trouve

\textbf{Scholie} des ateliers \add{(quasi-)} ensemblistes à figures de type fini : se donner un tel atelier, revient au même que de se donner un \struck{\ill{}} couple $(\mathcal{L}, \mathcal{M})$, $\mathcal{L}$ \ill{} ensemble, $\mathcal{M}$ ens. des figures élémentaires dans $\mathcal{L}$, telles que toute sous-figure $F$ d'un $X \in \mathcal{M}$ soit de type fini (\ill{} si les \ill{} \ill{}) et $X \in \mathcal{M}$ sont des figures finies) et satisfaisant de plus Atens$^{\sharp}$1 -- Atens$^{\sharp}$3. On reconstitue $\mathfrak{F}$ comme un sous-atelier de $\mathrm{Fig}(\mathcal{L})$, formé des figures \add{ens.} \struck{de type fini} \emph{de type fini} telles que ses multistrates soient dans $\mathcal{M}$.

Dans le cas d'un atelier ens. \add{\uncertain{fig.}} de type fini, on \struck{s'attendrait} \struck{\ill{} que pour une figure ensembliste $F$ dans $\mathcal{L}$, dont les multistrates sont dans $\mathcal{M}$, l'appartenance à $\mathfrak{F}$ ne dépend que de son \emph{ombre}.}\note{les quatre dernières lignes sont barrées de traits obliques ; la phrase est reprise en haut de la page 57}

\page{57}
on s'attend pas \add{\uncertain{pourtant}} que $\mathfrak{F}$ soit une sous-atelier spécieux, au sens \struck{le plus} fort où nous avons introduit ce terme, i.e. stable par Sup \emph{quelconques}. \struck{Ainsi}, on ne s'attend pas en général que \struck{\ill{}} $\lbrace F_x = \lbrace \lbrace x \rbrace \rbrace \rbrace_{x \in \mathcal{L}}$ soit dans $\mathfrak{F}$ ! La condition raisonnable \ill{} \ill{} que $\mathfrak{F}$ soit \ill{} \ill{} dès lors, c'est que $\mathfrak{F}$ soit \emph{finiment} spécieux. Vu la condition Atens$^{\sharp}$3, cela signifie simplement que la relation $\underset{\mathcal{M}}{\between}$ est la relation \struck{Atens$^{\sharp}$2} $\between$ ensembliste. Donc

\marginal{On dit alors que $\mathfrak{F}$ est un atelier ensembliste finiment spécieux}\note{note oblique dans la marge gauche}

\textbf{Corollaire du Scholie} Les ateliers (quasi) ensemblistes finiment spécieux, \add{\ill{}} de type fini, correspondent \ill{} figures de type fini
\[
  (\mathcal{L}, \mathcal{M}),
\]
$\mathcal{L}$ un ens. et $\mathcal{M}$ un ens. de figures élémentaires dans $\mathcal{L}$, satisfaisant Atens$^{\sharp}$1 et Atens$^{\sharp}$2 (i.e. les $\lbrace \lbrace x \rbrace \rbrace$ sont dans $\mathcal{M}$ et si $X \in \mathcal{M}$ \ill{} multistrates de $X$ \ill{}), \struck{tel} dans $\mathcal{M}$ \struck{tels} tel que $\forall X \in \mathcal{M}$ \struck{toute} sous-figure de $\mathcal{L}$ soit de t.f.\note{sic, « sous-figure de $\mathcal{L}$ », pour une sous-figure de $X$} Alors $\mathfrak{F}$ est l'ens. des

\page{58}
figures ensemblistes de type fini dont les multistrates sont dans $\mathcal{M}$.

\textbf{Remarque} En dehors du cas \add{finiment} spécieux, le seul \add{autre} cas intéressant que j'ai rencontré, est celui où la \struck{\ill{}} relation $\underset{\mathcal{M}}{\between}$ est celle-ci : $X$ et $Y$ ensemblistement compatibles, et $X \cap Y$ est \struck{\ill{}} \add{ou bien vide, ou bien} élémentaire i.e. est une multistrate commune.

Dans ce cas, $\mathfrak{F}$ a la propriété supplémentaire suivante. \struck{\ill{}}
\begin{itemize}
  \item[(1.57)] $\forall F \in \mathfrak{F}$, $\widetilde{F}$ admet le inf de \struck{deux} éléments \struck{i.e. tout} ou ce qui revient au même : si $X, Y \in \widetilde{F}$, \ill{} \ill{} \ill{} alors $X \cap Y$ (est l'inf pour $\leq$) est une multistrate.
\end{itemize}

\note{la page se termine par un double trait horizontal}

\page{59}
J'en viens à des \emph{axiomes de disjonction} (At 8 était l'« axiome des lieux »).\note{la première mouture (GF IV, dossier 156-4, p. 86) annonçait de même, sous le titre « Axiomes de disjonction et de support », qu'il était temps d'introduire les lieux ; elle s'arrête avant les axiomes de disjonction eux-mêmes}

Notations et d'abord
\struck{At 9 \ill{} Strict \ill{}}

\textbf{Proposition} Soit $\mathfrak{F}$ un atelier \add{\ill{}} $\vDash$ At 1 -- At 3 (on n'y va pas utiliser $\ll$). Si $F$ et $G \in \mathfrak{F}$, alors
\[
  F \parallel G \Longleftrightarrow \forall X \in \widetilde{F},\ Y \in \widetilde{G},\quad X \parallel Y .
\]
Immédiat. Également, \struck{\ill{}} on voudrait
\begin{itemize}
  \item[(1.58)] \struck{Corollaire} \struck{Proposition} Si $F \parallel G$, $F' \ll F$, $G' \ll G$, alors $F' \parallel G'$.
\end{itemize}
En effet, \struck{\ill{}} on est ramené au cas \add{(\ill{} \ill{})} $F = X'$, $G = Y'$, puis au cas $F = X$, $G = Y$, qu'on va utiliser At 9.

\marginal{Remarque \ill{} pour $X \parallel Y \Longrightarrow \mathrm{Omb}(X) \cap \mathrm{Omb}\, Y = \emptyset$ \ill{} \ill{} $\mathrm{omb}(X) \cap \mathrm{omb}(Y) \ldots = 0$}\note{note oblique dans la marge gauche, en haut, lue en partie}

Mais \struck{dans} dans le cas ensembliste à figures de t.f. (\ill{} \ill{} fini) cela ne semble pas résulter \ill{} \ill{} \ill{} \ill{} p. 55, \ill{} \ill{} Atens$^{\sharp}$3 \ill{} \ill{} dit quelles conditions \ill{} : la relation $\underset{\mathcal{M}}{\between}$ \struck{\ill{}} : d'être stable par passage aux sous-strates, et d'être plus forte que la relation $\underset{\mathrm{ens}}{\between}$.

\marginal{NB Pb, si $F \parallel G$, \ill{} \ill{} $Z \in \mathcal{M}$, $Z \ll F$, $Z \ll G$ \ill{} \ill{} ; $F \vee G$ \ill{} \ill{} \ill{} \ill{} donc \ill{} $\ll$ \ill{} \ill{}}\note{longue note oblique dans la marge gauche de la moitié inférieure, écrite serrée et en partie biffée, lue par fragments}

\page{60}
\uncertain{Ça n'a} pas l'air d'impliquer la condition. \ill{} \ill{} \ill{} pas que \struck{la disjonction} la disjonction ensembliste implique la disjonction géométrique. Prenons p.ex. $\mathcal{M}$ réduit : \struck{trois} \add{des} multistrates \struck{\ill{}} réduites à une seule \struck{ensemble} \add{strate} (ce qui rend vide la première condition Atens$^{\sharp}$3) et par exemple définie par les trois ensembles $A$, $B$ et $A'$
\[
  A' \subsetneq A, \qquad A \cap B = \emptyset .
\]
\struck{Prenons} \struck{donc}
\[
  X = \lbrace A \rbrace, \quad Y = \lbrace B \rbrace, \quad X' = \lbrace A' \rbrace, \quad \text{et soit } Y' = Y,
\]
\struck{On définit} Le \struck{donc}
\[
  X' \ll X, \qquad Y' = Y \text{ donc } Y' \ll Y .
\]
\note{le premier signe $\ll$ est surchargé}
donc \struck{\ill{}} \ill{} compatibilité ensembliste des trois \struck{$\between$ comme étant \ill{} la \ill{}} \struck{At} multistrates est déjà acquise. Donc on peut prendre pour $\underset{\mathcal{M}}{\between}$ n'importe quelle relation sym. refl., \struck{donc} \add{n'importe quelle} partie \struck{\ill{}} de l'ens. des trois paires de \struck{strates} multistrates $X, Y, X'$ --- on choisira
\[
  X \underset{\mathcal{M}}{\between} Y, \quad \text{et} \quad X' \underset{\mathcal{M}}{\overline{\between}} Y
\]
alors
\[
  X \underset{\mathfrak{F}}{\parallel} Y, \quad \text{mais} \quad X' \underset{\mathfrak{F}}{\overline{\parallel}} Y, \quad \text{ok.}
\]

On posera donc
\begin{itemize}
  \item[(*)] \struck{\ill{}} Si $X, Y$ \struck{\ill{}}, $X', Y' \in \mathcal{M}$, $X' \ll X$, $Y' \ll Y$, et si $X \parallel Y$, alors $X' \parallel Y'$ (il suffit de poser $X \between Y$)
\end{itemize}
\note{devant « Si », un mot fortement biffé, sans doute un nom d'axiome (« At \ldots ») ; « (*) » est écrit dans la marge. Le texte se poursuit dans le lot suivant}

\end{document}
